\chapter{Background}
\label{ch:background}

This chapter introduces the concepts on which the thesis builds. We first
review how multilingual capabilities arise in large language models (LLMs)
and where they break down (Section~\ref{sec:multilingual-lms}). We then
discuss the special role of English as a pivot language, both as an explicit
engineering device in classical multilingual NLP and as an implicit property
of the internal representations of English-centric LLMs
(Section~\ref{sec:english-pivot}). Section~\ref{sec:multilingual-eval}
surveys how multilingual ability is evaluated today and motivates
alignment-based, reference-free metrics such as MEXA. Finally,
Section~\ref{sec:parallel-corpora} describes the parallel corpora that make
such evaluations possible and the resource disparities that constrain them.

% ------------------------------------------------------------
\section{Multilingual Language Models}
\label{sec:multilingual-lms}

% ..............................................
\subsection{Architectures and Pre-training}
\label{subsec:architectures}

Virtually all modern language models are based on the Transformer
architecture \citep{vaswani2017attention}, a stack of identical layers that
combine multi-head self-attention with position-wise feed-forward networks.
Each layer transforms a sequence of token representations into a new
sequence of contextualized representations; the output of layer $l$ for
token $t$, the \emph{hidden state} $h_t^{(l)}$, is the central object of
study in this thesis, because it is from these layer-wise hidden states
that we extract sentence embeddings.

Two architectural families dominate multilingual NLP. \emph{Encoder-only}
models such as mBERT \citep{devlin2019bert} and XLM-R
\citep{conneau2020xlmr} use bidirectional attention and are pre-trained
with masked-language-modeling objectives on deliberately multilingual
corpora (XLM-R covers 100 languages sampled from CommonCrawl). Their
representations are designed to serve downstream classification and
retrieval tasks, and dedicated sentence encoders such as LaBSE
\citep{feng2022labse}, LASER3 \citep{heffernan2022laser3}, and multilingual
E5 \citep{wang2024multilingual} additionally fine-tune on parallel data
with contrastive or translation-ranking objectives so that translations map
to nearby points in embedding space. \emph{Decoder-only} models such as
GPT-3 \citep{brown2020gpt3} and the Llama family \citep{dubey2024llama3}
use causal (left-to-right) attention and are pre-trained autoregressively
to predict the next token. They are the dominant paradigm for generative
LLMs, but their pre-training corpora are heavily skewed toward English:
for Llama~3, roughly 90\% of the pre-training tokens are English
\citep{dubey2024llama3}. Models of this kind are commonly called
\emph{English-centric}, even though they exhibit non-trivial abilities in
dozens of languages.

This asymmetry creates the central question of the thesis: an
English-centric decoder model is never explicitly trained to align
languages, yet it acquires multilingual competence as a by-product of
incidental multilingual data in its corpus. Understanding \emph{where} and
\emph{how well} this competence materializes inside the network --- rather
than only measuring it behaviorally on downstream tasks --- is what
alignment-based evaluation aims to do.

% ..............................................
\subsection{The Tokenization Bottleneck}
\label{subsec:tokenization}

Before any text reaches the Transformer layers, it is segmented into
subword units by a tokenizer, typically trained with byte-pair encoding
\citep{sennrich2016bpe} or a unigram language model as implemented in
SentencePiece \citep{kudo2018sentencepiece}. The tokenizer's vocabulary is
learned on (a sample of) the pre-training corpus, and therefore inherits
its language distribution: frequent English words receive dedicated tokens,
while words in under-represented languages are shattered into many short
fragments, in the worst case into individual bytes.

The severity of this effect is commonly quantified by \emph{fertility},
the average number of subword tokens produced per word (or per character
for scripts without whitespace word boundaries). High fertility has three
practical consequences. First, it inflates sequence lengths, making
inference disproportionately expensive for low-resource languages. Second,
it dilutes the semantic content carried by each token: a model that sees
Amharic or Tibetan only as byte sequences must reconstruct lexical meaning
from much noisier units than it does for English. Third, and most relevant
to this thesis, it affects the quality of sentence embeddings built from
token hidden states, since averaging over many weakly informative fragments
yields blurrier sentence representations. Non-Latin scripts (Ge'ez,
Myanmar, Khmer, Cherokee) and low-resource languages regardless of script
are the most affected. Throughout our experiments we therefore track
tokenizer fertility on both evaluation corpora alongside the alignment
scores, which allows us to distinguish genuinely weak cross-lingual
representations from representations that are handicapped already at the
tokenization stage.

% ..............................................
\subsection{Cross-lingual Transfer}
\label{subsec:transfer}

The practical value of multilingual models rests on \emph{cross-lingual
transfer}: capabilities acquired predominantly in one language (in
practice, English) become available in other languages without
language-specific supervision. The standard explanation is that the model
learns a partially \emph{shared semantic space} in which translation
equivalents --- words or sentences with the same meaning --- receive
similar representations, so that task knowledge attached to a region of
that space is accessible from any language that maps into it
\citep{hammerl2024survey}.

Empirical evidence supports this picture with an important refinement:
the degree of sharing is \emph{layer-dependent}. Input embeddings and the
first layers remain largely language-specific, dominated by lexical and
orthographic features; representations become increasingly
language-neutral through the middle of the network; and the final layers
re-specialize toward the output vocabulary of the target language.
Interpretability studies on Llama-family models show precisely this
trajectory \citep{wendler2024llamas,zhao2024multilingualism}, and it
motivates the layer-wise design of MEXA: if cross-lingual alignment exists
anywhere in an English-centric model, it should be found in the
intermediate layers, and a single-layer (e.g., last-layer) probe would
systematically underestimate it.

\citet{hammerl2024survey} organize the methods that \emph{create} or
\emph{improve} such alignment into families --- joint multilingual
pre-training, alignment objectives on parallel data, contrastive
sentence-level training, and post-hoc mapping between monolingual spaces.
This thesis is concerned with the complementary problem: not producing
alignment, but \emph{measuring} the alignment that a given model already
has, layer by layer and language by language.

% ------------------------------------------------------------
\section{English as a Pivot Language}
\label{sec:english-pivot}

% ..............................................
\subsection{Pivot Translation and Alignment}
\label{subsec:pivot-translation}

Using English as a bridge between languages long predates neural LLMs. In
statistical machine translation, direct parallel data between two
non-English languages is often scarce, so translation is routed through a
\emph{pivot} language: source $\rightarrow$ English $\rightarrow$ target
\citep{wu2007pivot}. Because English participates in the largest number of
parallel corpora, it is the pivot of choice, and the quality of any
source--target system becomes bounded by the quality of the two
English-anchored legs. Neural machine translation reproduced the pattern
in a softer form: multilingual NMT systems trained on English-centric data
develop an ``interlingua-like'' internal representation that enables
zero-shot translation between language pairs never seen together in
training \citep{johnson2017zeroshot}.

Modern decoder-only LLMs internalize this bridge. Using the logit-lens
technique, \citet{wendler2024llamas} show that when Llama-2 models process
non-English input, the intermediate layers pass through a phase in which
the representations are closest to \emph{English} tokens before the final
layers rotate them toward the target-language vocabulary --- the model
demonstrably ``works in English'' in its middle layers.
\citet{zhao2024multilingualism} corroborate this for reasoning tasks,
finding that multilingual inputs are converted into English-centric latent
representations in which the actual computation takes place. This
\emph{pivot hypothesis} is the theoretical foundation of MEXA: if
non-English understanding is routed through English-like intermediate
representations, then the degree to which a language's representations
align with English is a direct measurement of how much of the model's
(English-acquired) capability that language can access.

% ..............................................
\subsection{English Centricity and Bias}
\label{subsec:english-bias}

The pivot mechanism is a double-edged sword. On one side, it is exactly
what makes cross-lingual transfer from a 90\%-English corpus possible at
all. On the other, it builds a structural asymmetry into the model.
Languages that map cleanly onto English-like representations inherit the
model's full capability; languages that do not --- because of typological
distance, script, or sheer data scarcity --- are effectively locked out,
regardless of how much task knowledge the model possesses. Concepts and
distinctions that English does not encode (grammatical evidentiality,
honorific systems, culture-specific terminology) risk being flattened when
they are forced through an English-shaped bottleneck, and cultural framing
in generated text tends to default to Anglophone norms.

For evaluation, English centricity raises a methodological question that
this thesis addresses empirically: is English \emph{privileged} as a
measurement anchor, or merely \emph{convenient}? MEXA, as proposed, always
measures alignment against English. If the pivot hypothesis is the true
mechanism, replacing English with another pivot language should
systematically lower the measured alignment, in proportion to how far the
alternative pivot is from the model's internal hub. Our non-English-pivot
experiments (Section~\ref{sec:experimental-settings}) test exactly this by
recomputing alignment against French, German, Arabic, Chinese, and Basque
pivots.

% ------------------------------------------------------------
\section{Multilingual Evaluation}
\label{sec:multilingual-eval}

% ..............................................
\subsection{Standard Benchmarks}
\label{subsec:benchmarks}

The most direct way to assess multilingual capability is to run the model
on parallel downstream tasks. FLORES-200 \citep{nllbteam2022} provides
professionally translated Wikimedia sentences in 204 language--script
pairs and is the de-facto standard for evaluating translation and
cross-lingual retrieval. Belebele \citep{bandarkar2024belebele} builds a
machine reading comprehension benchmark on top of FLORES passages,
covering 122 language variants with fully parallel multiple-choice
questions --- crucially, differences in scores across languages can only
stem from the language itself, not from differing content. Knowledge- and
reasoning-oriented English benchmarks such as MMLU
\citep{hendrycks2021mmlu} and ARC have been machine-translated into
26--31 languages to form m-MMLU and m-ARC \citep{lai2023okapi}.

These benchmarks are indispensable, but they have structural limits: they
exist for at most a couple of hundred languages (Belebele: 122, m-MMLU:
34, m-ARC: 31), they are expensive to extend because each new language
needs high-quality translated task data, and machine-translated variants
inherit translation noise. For the several thousand remaining languages,
task-based evaluation is simply unavailable --- which is the gap that
alignment-based estimation aims to fill. In this thesis, downstream
benchmarks play the role of \emph{validation targets}: MEXA scores are
correlated against Belebele and m-ARC performance to test whether
alignment is a faithful proxy for task capability.

% ..............................................
\subsection{Reference-Based vs.\ Reference-Free Metrics}
\label{subsec:metrics}

Evaluation metrics in multilingual NLP fall into two broad families.
\emph{Reference-based} metrics compare model output against gold
references. Surface-overlap metrics such as BLEU
\citep{papineni2002bleu} are cheap and language-agnostic in principle, but
they reward n-gram overlap rather than meaning, penalize legitimate
paraphrases, and interact badly with morphologically rich languages and
non-whitespace scripts. Learned metrics such as COMET \citep{rei2020comet}
correlate far better with human judgments by scoring
source--hypothesis--reference triples with a fine-tuned multilingual
encoder, but they require human references and a trained scoring model,
and their reliability degrades precisely for the low-resource languages
where evaluation is most needed.

\emph{Reference-free} (quality-estimation or model-based) metrics remove
the dependency on gold outputs. Embedding-similarity approaches compare
representations of source and translation directly
\citep{reimers2020multilingual}, and model-internal approaches --- the
category MEXA belongs to --- do not look at generated text at all.
Instead, they probe the model's hidden states on parallel \emph{inputs}
and quantify structure that is known to correlate with capability. The
advantages are substantial: no task data, no references, no additional
trained judge, and applicability to any language for which a hundred
parallel sentences exist. The corresponding risk is validity --- an
internal signal is only useful insofar as it predicts external behavior
--- which is why correlation with downstream benchmarks
(Section~\ref{subsec:benchmarks}) is an essential part of the methodology
rather than an afterthought.

% ..............................................
\subsection{Challenges in Cross-lingual Alignment Evaluation}
\label{subsec:alignment-challenges}

Measuring alignment from raw embedding geometry is less straightforward
than it appears, for reasons rooted in the geometry of learned
representation spaces.

\paragraph{Anisotropy.} Transformer embedding spaces are strongly
anisotropic: representations occupy a narrow cone rather than spreading
over the full space, so the cosine similarity between two \emph{unrelated}
sentences is already high \citep{rajaee2022isotropy}. In multilingual
models the effect is compounded by outlier dimensions with extreme
variance that dominate cosine computations; removing a handful of such
dimensions can change similarity rankings entirely
\citep{hammerl2023anisotropy}. Absolute cosine values are therefore not
comparable across models, layers, or languages.

\paragraph{Hubness.} In high-dimensional spaces, some points become
``hubs'' that appear among the nearest neighbors of disproportionately
many queries. A naive nearest-neighbor criterion over cosine similarities
rewards such hubs and distorts retrieval-based measurements.

\paragraph{Language identity vs.\ semantics.} Embeddings encode both
\emph{what} a sentence says and \emph{which language} it is written in.
Two sentences in the same language are often closer to each other than a
sentence and its translation, simply because of the shared language
signal. A useful alignment metric must isolate the semantic component
from this language-identity component.

MEXA's design responds to all three problems at once: instead of using
cosine similarity as an absolute quantity, it uses it only
\emph{comparatively}, asking whether the correct translation outranks all
distractors within the same similarity matrix (a bidirectional
precision-at-1 criterion, formalized in
Section~\ref{sec:mexa-pipeline}). Rank-based comparisons within a fixed
matrix are invariant to the global scale and offset distortions that
anisotropy introduces, and the bidirectional requirement suppresses hub
sentences that would pass a one-directional test. As an extension, this
thesis additionally evaluates a \emph{mean-centered} variant in which each
language's mean embedding is subtracted before computing similarities,
directly removing the language-identity component discussed above.

% ------------------------------------------------------------
\section{Parallel Corpora}
\label{sec:parallel-corpora}

% ..............................................
\subsection{Sources and Typology}
\label{subsec:corpora-sources}

Alignment-based evaluation requires \emph{parallel} data: the same
sentences translated across many languages, so that semantic content is
held constant while language varies. Parallel corpora differ along
several axes that matter for measurement quality: translation quality
(professional vs.\ crawled/mined, e.g.\ OPUS \citep{tiedemann2012opus}),
domain (news, web, religious text), degree of parallelism (bitexts vs.\
$n$-way parallel corpora), and language coverage.

Two corpora sit at opposite ends of the coverage--quality trade-off and
are both used in this thesis. \textbf{FLORES-200} \citep{nllbteam2022}
contains professionally translated, quality-assured sentences from
Wikimedia sources in 204 language--script pairs; it is $n$-way parallel,
modern-domain, and high quality, but its coverage stops at roughly 200
languages. The \textbf{Parallel Bible Corpus} \citep{mayer2014bible}
exploits the fact that the Bible is the most translated text in existence,
with verse-level alignment across more than 1{,}400 language--script
pairs; the derived \emph{super-parallel} subset (sPBC) restricts to verses
available in all covered languages. Its domain is narrow and archaic and
translation styles vary, but it is the only resource that reaches deep
into the low-resource tail, which is exactly where reference-free
evaluation is most valuable.

% ..............................................
\subsection{Resource Disparity}
\label{subsec:resource-disparity}

The distribution of parallel (and monolingual) data across the world's
languages is extremely skewed. A few dozen high-resource languages account
for the overwhelming majority of available text; thousands of languages
have effectively no digital footprint beyond a Bible translation. This
disparity propagates through the entire pipeline studied in this thesis:
low-resource languages receive less pre-training data (weaker
representations), fragmented tokenization
(Section~\ref{subsec:tokenization}), and no downstream benchmarks
(Section~\ref{subsec:benchmarks}) --- so their capability can be neither
learned as well nor measured as easily.

For evaluation design, resource disparity has two concrete consequences.
First, any claim about ``multilingual'' capability that is validated only
on high-resource languages is systematically optimistic; extending
measurement to the sPBC's low-resource tail is therefore a substantive
robustness test, not an ornament. Second, corpus effects must be
disentangled from language effects: a language may score lower on the
Bible corpus than on FLORES because the model is genuinely weaker in it,
or because archaic religious register is far from the model's training
distribution. Using both corpora on the same models, as done here, makes
this confound visible and quantifiable.


